\section{Quadrados de ordinais}\label{sec:quadradosDeOrdinais}
Vimos que $\omega\times \omega$ é infinito enumerável.

Tal fato não é específico de $\aleph_0$.
De fato, para cada cardinal infinito $\kappa$, $\kappa\times \kappa$ é equipotente a $\kappa$.

Para mostrar isso, utilizaremos algumas definições e resultados sobre ordinais.

\begin{definition}\label{def:ordemCanonicaEmOnQuadrado}
    Dizemos que um conjunto está em $\ON\times \ON$ se é um par ordenado de ordinais.

    Sejam $(\alpha, \beta)$ e $(\alpha', \beta')$ pares de ordinais.
    A ordem canônica em $\ON\times \ON$ é a relação $\triangleleft$ definida por:
    \begin{align*}
    (\alpha, \beta) \triangleleft (\alpha', \beta') \Longleftrightarrow &(\max\{\alpha, \beta\} < \max\{\alpha', \beta'\}) \\
    \lor &(\max\{\alpha, \beta\} = \max\{\alpha', \beta'\} \land \alpha < \alpha') \\
    \lor &(\max\{\alpha, \beta\} = \max\{\alpha', \beta'\} \land \alpha = \alpha' \land \beta < \beta').
    \end{align*}
\end{definition}
\begin{lemma}\label{lem:boaOrdemCanonicaNoQuadradoDeOrdinal}
    Para cada ordinal $\alpha$, $\alpha\times \alpha$ é uma boa ordem em $\ON\times \ON$, ou seja, para todos $(\alpha, \beta)$, $(\alpha', \beta')$ e $(\alpha'', \beta'')$ pares de ordinais, valem as seguintes propriedades.
    \begin{enumerate}[label=(\alph*)]
        \item Se $(\alpha, \beta) \triangleleft (\alpha', \beta')$ e $(\alpha', \beta') \triangleleft (\alpha'', \beta'')$, então $(\alpha, \beta) \triangleleft (\alpha'', \beta'')$.
        \item $(\alpha, \beta) \ntriangleleft (\alpha, \beta)$.
        \item $(\alpha, \beta) \triangleleft (\alpha', \beta') \lor (\alpha, \beta) = (\alpha', \beta') \lor (\alpha', \beta') \triangleleft (\alpha, \beta)$.
        \item Se $X$ é um conjunto não vazio de pares de ordinais, então $X$ tem um $\triangleleft$-menor elemento.
    \end{enumerate}
\end{lemma}

\begin{proof}
    (a) Suponha que $(\alpha, \beta) \triangleleft (\alpha', \beta')$ e $(\alpha', \beta') \triangleleft (\alpha'', \beta'')$.
    
    Note que, em qualquer caso, $\max\{\alpha, \beta\} \leq \max\{\alpha', \beta'\} \leq \max\{\alpha'', \beta''\}$.
    Se alguma das desigualdades forem estritas, então $(\alpha, \beta) \triangleleft (\alpha'', \beta'')$.
    Caso contrário, $\max\{\alpha, \beta\} = \max\{\alpha', \beta'\} = \max\{\alpha'', \beta''\}$.
    
    Nesse caso, temos que $\alpha \leq \alpha' \leq \alpha''$.
    Se alguma das desigualdades forem estritas, então $(\alpha, \beta) \triangleleft (\alpha'', \beta'')$.
    Caso contrário, $\alpha = \alpha' = \alpha''$.

    Nesse caso, temos que $\beta < \beta' < \beta''$, de modo que $(\alpha, \beta) \triangleleft (\alpha'', \beta'')$.

    (b) é imediato da definição.

    (c) Suponha que $(\alpha, \beta) \ntriangleleft (\alpha', \beta')$ e $(\alpha', \beta') \ntriangleleft (\alpha, \beta)$.
    Então, $\max\{\alpha, \beta\} = \max\{\alpha', \beta'\}$, $\alpha = \alpha'$, e $\beta = \beta'$.
    Assim, $(\alpha, \beta) = (\alpha', \beta')$.

    (d) Seja $X$ um conjunto não vazio de pares de ordinais.
    Seja $X_0=\{\max\{\alpha, \beta\} : (\alpha, \beta)\in X\}$.
    Note que $X_0$ é um conjunto de ordinais, e, portanto, tem um menor elemento $\xi$.

    Seja $Y=\{(\alpha, \beta)\in X : \max\{\alpha, \beta\} = \xi\}$.
    Seja $Y_0=\{\alpha : (\alpha, \beta)\in Y\}$.
    Note que $Y_0$ é um conjunto de ordinais, e, portanto, tem um menor elemento $\alpha_0$.

    Seja $Z=\{(\alpha, \beta)\in Y : \alpha = \alpha_0\}$.
    Seja $Z_0=\{\beta : (\alpha, \beta)\in Z\}$.
    Note que $Z_0$ é um conjunto de ordinais, e, portanto, tem um menor elemento $\beta_0$.

    Temos que $(\alpha_0, \beta_0)\in Z \subseteq Y \subseteq X$.

    Se $(\alpha, \beta)\in X$, temos que $\max\{\alpha, \beta\} \geq \xi$.
    Se $\max\{\alpha, \beta\} > \xi$, então $(\alpha_0, \beta_0) \triangleleft (\alpha, \beta)$;
    caso contrário, $\max\{\alpha, \beta\} = \xi$, e, assim, $(\alpha, \beta)\in Y$.

    Como $(\alpha, \beta)\in Y$, temos que $\alpha \geq \alpha_0$.
    Se $\alpha > \alpha_0$, então $(\alpha_0, \beta_0) \triangleleft (\alpha, \beta)$;
    caso contrário, $\alpha = \alpha_0$, e, assim, $(\alpha, \beta)\in Z$.

    Como $(\alpha, \beta)\in Z$, temos que $\beta \geq \beta_0$.
    Dessa forma, $(\alpha_0, \beta_0) \trianglelefteq (\alpha, \beta)$.
\end{proof}

Se $(\alpha, \beta)\in \ON\times \ON$, denotaremos, nesta seção, $(\alpha, \beta)_{\downarrow}=\{(\alpha', \beta')\in \ON\times \ON : (\alpha', \beta') \triangleleft (\alpha, \beta)\}$, que é um conjunto pois $(\alpha, \beta)_{\downarrow} \subseteq \xi\times \xi$, onde $\xi$ é qualquer ordinal com $\xi>\max\{\alpha, \beta\}$.

Note que $(\alpha, 0)_{\downarrow} = \alpha\times \alpha$ para cada ordinal $\alpha$.

\begin{theorem}\label{thm:quadradoDeCardinalInfinito}
    Para cada cardinal infinito $\kappa$, $\tp(\kappa\times \kappa, \triangleleft) = \kappa$.
\end{theorem}
\begin{proof}
    Procedemos por indução transfinita para mostrar que para todo ordinal $\alpha$, $\tp(\aleph_\alpha\times \aleph_\alpha, \triangleleft) = \aleph_\alpha$.

    Como $\aleph_\alpha\preceq \aleph_\alpha\times \aleph_\alpha$ para cada ordinal $\alpha$, temos que $\tp(\aleph_\alpha\times \aleph_\alpha, \triangleleft) \geq |\aleph_\alpha\times \aleph_\alpha| \geq |\aleph_\alpha| = \aleph_\alpha$.
    Assim, resta mostrar que, para cada ordinal $\alpha$, $\tp(\aleph_\alpha\times \aleph_\alpha, \triangleleft) \leq \aleph_\alpha$.

    Para $\alpha=0$, temos que, se $(n, m)\in \omega\times \omega$, então $(n, m)_{\downarrow} \subseteq \max\{n, m\}\times \max\{n, m\}$, que é finito.
    Assim, $(\omega\times \omega, \triangleleft)$ é uma boa ordem infinita tal que todo segmento inicial próprio é finito, de modo que $\tp(\omega\times \omega, \triangleleft) = \omega$.

    Suponha agora que $\alpha>0$ e suponha por absurdo que $\tp(\aleph_\alpha\times \aleph_\alpha, \triangleleft) > \aleph_\alpha$.
    Seja $\xi=\tp(\aleph_\alpha\times \aleph_\alpha, \triangleleft)$.
    Considere $f:\aleph_\alpha\times \aleph_\alpha\to \xi$ isomorfismo.

    Existem $\alpha_0, \beta_0 < \aleph_\alpha$ tais que $f(\alpha_0, \beta_0) = \aleph_\alpha$.
    Assim, $|(\alpha_0, \beta_0)_{\downarrow}|=\aleph_\alpha$.
    Se ambos $\alpha_0, \beta_0$ são finitos, temos que $|(\alpha_0, \beta_0)_{\downarrow}|\leq |\max\{\alpha_0+1, \beta_0+1\}\times \max\{\alpha_0+1, \beta_0+1\}|$ é infinito, um absurdo.

    Assim, $\gamma=\max\{\alpha_0, \beta_0\}+1$, é infinito.
    Temos que $|(\alpha_0, \beta_0)_{\downarrow}|\leq |\gamma\times \gamma|$.

    Seja $\delta$ tal que $|\gamma|=\aleph_\delta$.
    Então $|\gamma\times \gamma| \leq |\aleph_\delta\times \aleph_\delta|=\aleph_\delta$ por hipótese de indução.
    Assim, temos um absurdo:

    \[
        \aleph_\alpha = |(\alpha_0, \beta_0)_{\downarrow}| \leq |\gamma\times \gamma| \leq \aleph_\delta < \aleph_\alpha.
    \]
\end{proof}